% Part: proof-theory
% Chapter: sequent-calculus
% Section: rules-G3c

\documentclass[../../../include/open-logic-section]{subfiles}

\begin{document}

\begin{table}
\begin{tabular}{@{}c@{}c@{}}
\multicolumn{2}{@{}c@{}}{అక్షయాలు}\\[2ex]
$!C, \Gamma \Sequent \Delta, !C$ & $\lfalse, \Gamma \Sequent \Delta$
\\[2ex]
\multicolumn{2}{@{}c@{}}{తార్కిక నియమాలు}\\[2ex]
\Axiom$\lnot !A, \Gamma \fCenter \Delta, !A $
\RightLabel{\LeftR{\lnot}}
\UnaryInf$\lnot !A, \Gamma \fCenter \Delta$
\DisplayProof
&
\Axiom$!A, \Gamma \fCenter$
\RightLabel{\RightR{\lnot}}
\UnaryInf$ \Gamma \fCenter \Delta, \lnot !A $
\DisplayProof
\\[4ex]
\Axiom$ !A, !B, \Gamma \fCenter \Delta$
\RightLabel{\LeftR{\land}}
\UnaryInf$ !A \land !B, \Gamma \fCenter \Delta$
\DisplayProof
&
\Axiom$\Gamma \fCenter \Delta, !A$
\Axiom$ \Gamma \fCenter \Delta, !B$
\RightLabel{\RightR{\land}}
\BinaryInf$ \Gamma \fCenter \Delta, !A \land !B $
\DisplayProof
\\[4ex]\Axiom$!A, \Gamma \fCenter \Delta$
\Axiom$!B, \Gamma \fCenter \Delta$
\RightLabel{\LeftR{\lor}}
\BinaryInf$!A \lor !B, \Gamma \fCenter \Delta$
\DisplayProof
&
\Axiom$\Gamma \fCenter \Delta, !A, !B$
\RightLabel{\RightR{\lor}}
\UnaryInf$ \Gamma \fCenter \Delta, !A \lor !B$
\DisplayProof
\\[4ex]
\Axiom$!A \lif !B, \Gamma \fCenter \Delta, !A$
\Axiom$ !B, \Gamma \fCenter \Delta$
\RightLabel{\LeftR{\lif}}
\BinaryInf$!A \lif !B, \Gamma \fCenter \Delta$
\DisplayProof
&
\Axiom$ !A, \Gamma \fCenter !B$
\RightLabel{\RightR{\lif}}
\UnaryInf$ \Gamma \fCenter \Delta, !A \lif !B $
\DisplayProof
\\[4ex]
\Axiom$ !A(t), \lforall[x][!A(x)], \Gamma \fCenter \Delta$
\RightLabel{\LeftR{\lforall}}
\UnaryInf$ \lforall[x][!A(x)],\Gamma \fCenter \Delta$
\DisplayProof
&
\Axiom$ \Gamma \fCenter !A(a) $
\RightLabel{\RightR{\lforall}}
\UnaryInf$ \Gamma \fCenter \lforall[x][!A(x)]$
\DisplayProof
\\[4ex]
\Axiom$ !A(a), \Gamma \fCenter \Delta $
\RightLabel{\LeftR{\lexists}}
\UnaryInf$ \lexists[x][!A(x)], \Gamma \fCenter \Delta$
\DisplayProof
&
\Axiom$ \Gamma \fCenter \Delta, \lexists[x][!A(x)], !A(t) $
\RightLabel{\RightR{\lexists}}
\UnaryInf$ \Gamma \fCenter \Delta, \lexists[x][!A(x)]$
\DisplayProof
\end{tabular}
\caption{అంతఃప్రజ్ఞావాద తర్కానికి బహుళ-నిర్ధారణల \Log{mG3i} నియమాలు. $\RightR{\forall}$, $\LeftR{\exists}$లలో \tetoken{స్థిరాంకం}{!!{constant}}~$a$ నిర్ధారణ సీక్వెంట్‌లో రాకూడదు. $\Gamma$, $\Delta$ \tetoken{సూత్రాల}{!!{formula}s} బహుసమితులు. అక్షయాల్లో \tetoken{సూత్రం}{!!{formula}}~$!C$ పరమాణువు కావాలి. మూల పక్క గమనిక \Log{mG1m}ను, అక్షయాలు $\lfalse,
\Gamma \Sequent \Delta$ లేని \Log{mG1i}గా పేర్కొంటుంది.}
\ollabel{tab:mG3i}
\end{table}
\sourcecorrection{OLTEPTSEQMG3I-001}{సార్వత్రిక ఎడమ ఆధారంలో ప్రధాన సూత్రం, పక్క సందర్భం మధ్య కామాను చేర్చాం.}
\sourcecorrection{OLTEPTSEQMG3I-002}{ఈ పట్టిక బహుళ-నిర్ధారణల మూడవ వ్యవస్థను ఇస్తున్నా మూల పక్క గమనిక మొదటి వ్యవస్థ పేర్లను వాడింది; ఆ పేర్లకు వేరు నియమ నిర్వచనం మూల నిరూపణ-సిద్ధాంత పాఠ్యంలో కనిపించలేదు. మూల సూచనగా మాత్రమే నిలిపాం; ఇదే కనిష్ఠ తర్కానికి పూర్తి నిర్వచనమని ధృవీకరించలేదు.}
\sourcecorrection{OLTEPTSEQMG3I-003}{మూల పట్టిక గుర్తు సాంప్రదాయ వ్యవస్థ గుర్తునే మళ్లీ వాడింది; బహుళ-నిర్ధారణల పట్టికకు ప్రత్యేక గుర్తు ఇచ్చాం.}

\end{document}
